% GENERATED FILE. Edit canonical lessons, then run npm run generate:books.
% canonical-source-hash: fnv1a64:dae75c57bb963215
% canonical-lessons: AR-C487-ghamara, AR-C487-jarafa, AR-C487-ihtama, AR-C487-lajaa, AR-C487-malja, AR-R487-review-487-487, AR-R487-review-final-487

\chapter{Floods and shelter}
\label{ch:ar-floods-and-shelter}
\hlchaptermodality{487}

\begin{chapteropening}
\textbf{By the end of this chapter:} \emph{I can say to submerge, to sweep away, to take shelter, to seek refuge and a shelter.}

Complete the last lesson of chapter 487: I can say to submerge, to sweep away, to take shelter, to seek refuge and a shelter.
\end{chapteropening}

\section[\texorpdfstring{\ar{غمر} (ghamara)}{غمر (ghamara)}]{\ar{غمر} (ghamara) — to submerge; to flood}

\label{lesson:AR-C487-ghamara}



\begin{quote}
Before the new one: say the Arabic for refreshing, then the Arabic for to be exposed.
\end{quote}

\subsection*{You'll want to know: \ar{غمر}}
\textbf{\ar{غمر}} — \emph{ghamara} — \textquotedblleft{}to submerge; to flood\textquotedblright{}.

\begin{quote}
\ar{غمر} \ar{الماء} \ar{الشارع} \ar{بعد} \ar{المطر}.
\end{quote}

\begin{quote}
\emph{ghamara al-māʾ ash-shāriʿ baʿda al-maṭar.}
\end{quote}

\begin{quote}
Water flooded the street after the rain.
\end{quote}

Present: \ar{يغمر}, \emph{yaghmur}.

\begin{cousinweb}
\textbf{\ar{غمر}} (\emph{ghamara}) is built on the native Arabic root \emph{gh-m-r}, \textquotedblleft{}abundant water; to cover\textquotedblright{} (Lane). It is the plain verb.
\end{cousinweb}

\begin{grammarlens}[title={What you've built}]
One more word for getting around.
\end{grammarlens}

\subsection*{Your turn}
\begin{itemize}
\raggedright
  \item \emph{Say it:} \emph{ghamara}
  \item \emph{Say it:} \emph{ghamara}, once more
  \item \emph{Say it:} \emph{ghamara al-māʾ ash-shāriʿ baʿda al-maṭar.} — Water flooded the street after the rain.
  \item \emph{Recall:} say the Arabic for refreshing, then the Arabic for to be exposed, then say \emph{ghamara} again
\end{itemize}

\begin{tcolorbox}[breakable,colback=teal!4,colframe=teal!35!black,title={Before you move on}]
What does \ar{غمر} mean? (\textquotedblleft{}To submerge; to flood\textquotedblright{}.) Say it once more.
\end{tcolorbox}
\section[\texorpdfstring{\ar{جرف} (jarafa)}{جرف (jarafa)}]{\ar{جرف} (jarafa) — to sweep away}

\label{lesson:AR-C487-jarafa}



\begin{quote}
Before the new one: say the Arabic for to be exposed, then the Arabic for to submerge.
\end{quote}

\subsection*{You'll want to know: \ar{جرف}}
\textbf{\ar{جرف}} — \emph{jarafa} — \textquotedblleft{}to sweep away\textquotedblright{}.

\begin{quote}
\ar{جرف} \ar{الماء} \ar{السيارات} \ar{في} \ar{الوادي}.
\end{quote}

\begin{quote}
\emph{jarafa al-māʾ as-sayyārāt fī al-wādī.}
\end{quote}

\begin{quote}
The water swept the cars away in the valley.
\end{quote}

Present: \ar{يجرف}, \emph{yajruf}.

The noun \emph{jurf} is a bank undercut by water, a cliff.

\begin{cousinweb}
\textbf{\ar{جرف}} (\emph{jarafa}) is built on the native Arabic root \emph{j-r-f}, \textquotedblleft{}to sweep away\textquotedblright{} (Lane). EtymDB lists Syriac relatives.
\end{cousinweb}

\begin{grammarlens}[title={What you've built}]
One more word for getting around.
\end{grammarlens}

\subsection*{Your turn}
\begin{itemize}
\raggedright
  \item \emph{Say it:} \emph{jarafa}
  \item \emph{Say it:} \emph{jarafa}, once more
  \item \emph{Say it:} \emph{jarafa al-māʾ as-sayyārāt fī al-wādī.} — The water swept the cars away in the valley.
  \item \emph{Recall:} say the Arabic for to be exposed, then the Arabic for to submerge, then say \emph{jarafa} again
\end{itemize}

\begin{tcolorbox}[breakable,colback=teal!4,colframe=teal!35!black,title={Before you move on}]
What does \ar{جرف} mean? (\textquotedblleft{}To sweep away\textquotedblright{}.) Say it once more.
\end{tcolorbox}
\section[\texorpdfstring{\ar{احتمى} (iḥtamā)}{احتمى (iḥtamā)}]{\ar{احتمى} (iḥtamā) — to take shelter}

\label{lesson:AR-C487-ihtama}



\begin{quote}
Before the new one: say the Arabic for to submerge, then the Arabic for to sweep away.
\end{quote}

\subsection*{You'll want to know: \ar{احتمى}}
\textbf{\ar{احتمى}} — \emph{iḥtamā} — \textquotedblleft{}to take shelter\textquotedblright{}.

\begin{quote}
\ar{احتمينا} \ar{من} \ar{المطر} \ar{تحت} \ar{الجسر}.
\end{quote}

\begin{quote}
\emph{iḥtamaynā min al-maṭar taḥta al-jisr.}
\end{quote}

\begin{quote}
We took shelter from the rain under the bridge.
\end{quote}

Present: \ar{يحتمي}, \emph{yaḥtamī}.

\begin{cousinweb}
\textbf{\ar{احتمى}} (\emph{iḥtamā}) is built on the native Arabic root \emph{ḥ-m-y}, \textquotedblleft{}to protect, to defend\textquotedblright{} (Lane). It is Form VIII, \textquotedblleft{}to protect oneself\textquotedblright{}.
\end{cousinweb}

\begin{grammarlens}[title={What you've built}]
One more word for getting around.
\end{grammarlens}

\subsection*{Your turn}
\begin{itemize}
\raggedright
  \item \emph{Say it:} \emph{iḥtamā}
  \item \emph{Say it:} \emph{iḥtamā}, once more
  \item \emph{Say it:} \emph{iḥtamaynā min al-maṭar taḥta al-jisr.} — We took shelter from the rain under the bridge.
  \item \emph{Recall:} say the Arabic for to submerge, then the Arabic for to sweep away, then say \emph{iḥtamā} again
\end{itemize}

\begin{tcolorbox}[breakable,colback=teal!4,colframe=teal!35!black,title={Before you move on}]
What does \ar{احتمى} mean? (\textquotedblleft{}To take shelter\textquotedblright{}.) Say it once more.
\end{tcolorbox}
\section[\texorpdfstring{\ar{لجأ} (lajaʾa)}{لجأ (lajaʾa)}]{\ar{لجأ} (lajaʾa) — to seek refuge}

\label{lesson:AR-C487-lajaa}



\begin{quote}
Before the new one: say the Arabic for to sweep away, then the Arabic for to take shelter.
\end{quote}

\subsection*{You'll want to know: \ar{لجأ}}
\textbf{\ar{لجأ}} — \emph{lajaʾa} — \textquotedblleft{}to seek refuge\textquotedblright{}.

\begin{quote}
\ar{لجأ} \ar{الناس} \ar{إلى} \ar{المدرسة} \ar{أثناء} \ar{الفيضان}.
\end{quote}

\begin{quote}
\emph{lajaʾa an-nās ilā al-madrasa athnāʾa al-fayaḍān.}
\end{quote}

\begin{quote}
People took refuge in the school during the flood.
\end{quote}

Present: \ar{يلجأ}, \emph{yaljaʾ}.

It takes \emph{ilā} before the place of refuge.

\begin{cousinweb}
\textbf{\ar{لجأ}} (\emph{lajaʾa}) is built on the native Arabic root \emph{l-j-ʾ}, \textquotedblleft{}to take refuge\textquotedblright{} (Lane). It is the plain verb.
\end{cousinweb}

\begin{grammarlens}[title={What you've built}]
One more word for getting around.
\end{grammarlens}

\subsection*{Your turn}
\begin{itemize}
\raggedright
  \item \emph{Say it:} \emph{lajaʾa}
  \item \emph{Say it:} \emph{lajaʾa}, once more
  \item \emph{Say it:} \emph{lajaʾa an-nās ilā al-madrasa athnāʾa al-fayaḍān.} — People took refuge in the school during the flood.
  \item \emph{Recall:} say the Arabic for to sweep away, then the Arabic for to take shelter, then say \emph{lajaʾa} again
\end{itemize}

\begin{tcolorbox}[breakable,colback=teal!4,colframe=teal!35!black,title={Before you move on}]
What does \ar{لجأ} mean? (\textquotedblleft{}To seek refuge\textquotedblright{}.) Say it once more.
\end{tcolorbox}
\section[\texorpdfstring{\ar{ملجأ} (maljaʾ)}{ملجأ (maljaʾ)}]{\ar{ملجأ} (maljaʾ) — a shelter; a refuge}

\label{lesson:AR-C487-malja}



\begin{quote}
Before the new one: say the Arabic for to take shelter, then the Arabic for to seek refuge.
\end{quote}

\subsection*{You'll want to know: \ar{ملجأ}}
\textbf{\ar{ملجأ}} — \emph{maljaʾ} — \textquotedblleft{}a shelter; a refuge\textquotedblright{}.

\begin{quote}
\ar{نام} \ar{الأطفال} \ar{في} \ar{الملجأ} \ar{حتى} \ar{الصباح}.
\end{quote}

\begin{quote}
\emph{nāma al-aṭfāl fī al-maljaʾ ḥattā aṣ-ṣabāḥ.}
\end{quote}

\begin{quote}
The children slept in the shelter until morning.
\end{quote}

\begin{cousinweb}
\textbf{\ar{ملجأ}} (\emph{maljaʾ}) is built on the native Arabic root \emph{l-j-ʾ}, \textquotedblleft{}to take refuge\textquotedblright{} (Lane). The pattern \emph{mafʿal} makes it the place of refuge.
\end{cousinweb}

\begin{grammarlens}[title={What you've built}]
One more word for getting around.
\end{grammarlens}

\subsection*{Your turn}
\begin{itemize}
\raggedright
  \item \emph{Say it:} \emph{maljaʾ}
  \item \emph{Say it:} \emph{maljaʾ}, once more
  \item \emph{Say it:} \emph{nāma al-aṭfāl fī al-maljaʾ ḥattā aṣ-ṣabāḥ.} — The children slept in the shelter until morning.
  \item \emph{Recall:} say the Arabic for to take shelter, then the Arabic for to seek refuge, then say \emph{maljaʾ} again
\end{itemize}

\begin{tcolorbox}[breakable,colback=teal!4,colframe=teal!35!black,title={Before you move on}]
What does \ar{ملجأ} mean? (\textquotedblleft{}A shelter; a refuge\textquotedblright{}.) Say it once more.
\end{tcolorbox}
\section[(dialogue)]{Weather words — a review}

\label{lesson:AR-R487-review-487-487}



\begin{quote}
Say the Arabic for to submerge and to sweep away.
\end{quote}

\subsection*{Your turn}
\begin{itemize}
\raggedright
  \item \emph{Say it:} to submerge, to sweep away, to take shelter, to seek refuge, a shelter
\end{itemize}

\begin{tcolorbox}[breakable,colback=teal!4,colframe=teal!35!black,title={Before you move on}]
To submerge? (\textbf{\ar{غمر}}, \emph{ghamara}.) To sweep away? (\textbf{\ar{جرف}}, \emph{jarafa}.) To take shelter? (\textbf{\ar{احتمى}}, \emph{iḥtamā}.) To seek refuge? (\textbf{\ar{لجأ}}, \emph{lajaʾa}.) A shelter? (\textbf{\ar{ملجأ}}, \emph{maljaʾ}.)
\end{tcolorbox}
\section[(dialogue)]{Floods and shelter — once more}

\label{lesson:AR-R487-review-final-487}



\begin{quote}
Say the Arabic for to submerge and to sweep away.
\end{quote}

\subsection*{Your turn}
\begin{itemize}
\raggedright
  \item \emph{Say it:} to submerge, to sweep away, to take shelter, to seek refuge, a shelter
\end{itemize}

\begin{tcolorbox}[breakable,colback=teal!4,colframe=teal!35!black,title={Before you move on}]
To submerge? (\textbf{\ar{غمر}}, \emph{ghamara}.) To sweep away? (\textbf{\ar{جرف}}, \emph{jarafa}.) To take shelter? (\textbf{\ar{احتمى}}, \emph{iḥtamā}.) To seek refuge? (\textbf{\ar{لجأ}}, \emph{lajaʾa}.) A shelter? (\textbf{\ar{ملجأ}}, \emph{maljaʾ}.)
\end{tcolorbox}
